% =====================================================================
%  theory/governing_equations.tex
%  Phase 1, Blueprint V5.0, Task C
%
%  Derivation of the bulk piezothermoelastic system from ONE thermodynamic
%  potential, with intermediate steps, sign conventions, stability conditions
%  and a dimensional audit.
%
%  SOURCE-ONLY: no TeX engine is installed in this environment. This file has
%  been inspected as text and its symbolic content is verified by
%  code/foundation_checks.py (checks B1 and B2). It has NOT been compiled.
% =====================================================================
\documentclass[11pt]{article}
\usepackage{amsmath,amssymb,amsthm,bm}
\usepackage{geometry}
\geometry{margin=1in}
\newcommand{\s}{\sigma}
\newcommand{\eps}{\varepsilon}
\newcommand{\bth}{\vartheta}
\newcommand{\D}{\bm{D}}
\newcommand{\E}{\bm{E}}
\newcommand{\bu}{\bm{u}}
\newcommand{\bq}{\bm{q}}
\theoremstyle{plain}
\newtheorem{remark}{Remark}
\newtheorem{proposition}{Proposition}

\title{Bulk piezothermoelastic formulation from a single thermodynamic potential\\
\large Theory note 1 of 4 --- Phase 1, Blueprint V5.0}
\author{Derived for the Piezothermoelastic Anti-Reflection programme}
\date{4 November 2026}

\begin{document}
\maketitle

\begin{abstract}
This note derives the coupled piezothermoelastic field equations --- kinematics,
constitutive relations, momentum balance, Gauss law and the Lord--Shulman heat
equation --- from a single quadratic Helmholtz free-energy density. It is
deliberately written as a chain of small steps so that every sign can be checked
against its generating term. Two structural points are made that the working
blueprint lists as requiring independent verification: (i) the operator
$(1+\tau_q\partial_t)$ in the Lord--Shulman equation multiplies the \emph{whole}
bracket, including the pyroelectric term, and this is a consequence of a single
Cattaneo flux law rather than a modelling choice; (ii) the pyroelectric
coefficient enters $\D$ and the entropy with opposite signs, and it is exactly
this that closes the Maxwell cross-relations. Every relation below is labelled;
checks B1 and B2 of \texttt{code/foundation\_checks.py} verify them symbolically.
\end{abstract}

\section{Scope, notation and status labels}
\label{sec:scope}

\paragraph{Geometry.} Cartesian $(x,y,z)$; the sagittal plane is $x$--$z$;
the material is homogeneous in $y$, so $\partial_y \equiv 0$ (plane strain in the
sense $\eps_{yy}=0$, $\eps_{xy}=\eps_{yz}=0$). The poling axis is $\bm{e}_3$,
which is taken to lie \emph{in} the sagittal plane (it is $z$).

\paragraph{Unknowns.} Four scalar fields,
\begin{equation}
  \bm{y} = (u_x,\; u_z,\; \varphi,\; \bth)^{\mathsf T},
  \label{eq:unknowns}
\end{equation}
comprising two mechanical components, one electric potential and one temperature
perturbation. These are the four fields of Blueprint Eq.~(1). Three physical
categories (mechanical, electrical, thermal) are carried by four scalars; the
formulation is not ``four independent physics''.

\paragraph{Conventions.}
\begin{itemize}
  \item Time dependence $e^{-i\omega t}$. Consequently $\partial_t \to -i\omega$
        everywhere, and a field behaving as $e^{i(\kappa x+\xi z-\omega t)}$
        \emph{decays} in $+z$ when $\operatorname{Im}\xi>0$.
  \item Small perturbations about a uniform reference temperature $T_0>0$:
        $\bth \equiv T - T_0$, $|\bth|/T_0 \ll 1$.
  \item Electric field $\E = -\nabla\varphi$, so $E_i = -\varphi_{,i}$.
  \item Small-strain tensor $\eps_{ij} = \tfrac12(u_{i,j}+u_{j,i})$.
  \item Voigt engineering shear $\gamma_{xz} = 2\eps_{xz} = u_{x,z}+u_{z,x}$;
        $\s_{xz}$ is conjugate to $\gamma_{xz}$, not to $\eps_{xz}$.
  \item Symmetry: transverse isotropy of class $6mm$ about $\bm{e}_3$, except in
        \S\ref{sec:isotropic} where the isotropic reduction is used.
\end{itemize}

\paragraph{Status labels.} \textsc{Derived}: obtained here with the shown steps.
\textsc{Tested}: additionally verified symbolically or numerically by
\texttt{code/foundation\_checks.py}. \textsc{Assumed}: asserted without proof.
\textsc{Unresolved}: open.

\section{Kinematics}
\label{sec:kinematics}

With $\partial_y = 0$ the only non-zero strain components are
\begin{equation}
  \eps_{xx} = u_{x,x}, \qquad
  \eps_{zz} = u_{z,z}, \qquad
  \gamma_{xz} = u_{x,z} + u_{z,x}.
  \label{eq:strain}
\end{equation}
For a partial mode $\bm{y} = \bm{a}\,e^{i(\kappa x + \xi z - \omega t)}$ with
$\bm{a}=(U,W,\Phi,\Theta)^{\mathsf T}$, the operators become
$\partial_x \to i\kappa$, $\partial_z \to i\xi$, and
\begin{equation}
  \eps_{xx} = i\kappa U, \qquad
  \eps_{zz} = i\xi W, \qquad
  \gamma_{xz} = i(\xi U + \kappa W),
  \label{eq:strain-mode}
\end{equation}
\begin{equation}
  E_x = -i\kappa\Phi, \qquad E_z = -i\xi\Phi.
  \label{eq:E-mode}
\end{equation}

\section{The generating potential}
\label{sec:potential}

\subsection{Construction}

Take the quadratic Helmholtz free-energy density per unit volume, expanded about
$(T_0,\ \eps=0,\ \E=0)$. For $6mm$ symmetry with the poling axis along $z$ and
plane strain, the independent constants are $C_{11},C_{12},C_{13},C_{33},C_{44}$
(elastic, at constant $\E$ and $\bth$), $e_{31},e_{33},e_{15}$ (piezoelectric),
$\eps_{11},\eps_{33}$ (permittivity at constant strain), $\beta_1,\beta_3$
(thermoelastic coupling) and $p$ (pyroelectric):
\begin{align}
  \Psi(\eps,\E,\bth) &=
  \tfrac12\Big[ C_{11}(\eps_{xx}^2+\eps_{yy}^2) + C_{33}\eps_{zz}^2
    + 2C_{13}(\eps_{xx}+\eps_{yy})\eps_{zz}
    + 2C_{12}\eps_{xx}\eps_{yy} + C_{44}\gamma_{xz}^2 \Big] \notag\\[2pt]
  &\quad - e_{31}(\eps_{xx}+\eps_{yy})E_z - e_{33}\eps_{zz}E_z
    - e_{15}\gamma_{xz}E_x \notag\\[2pt]
  &\quad - \tfrac12\big(\eps_{11}E_x^2 + \eps_{33}E_z^2\big) \notag\\[2pt]
  &\quad - \beta_1\bth(\eps_{xx}+\eps_{yy}) - \beta_3\bth\eps_{zz}
    - p\,\bth E_z \notag\\[2pt]
  &\quad - \tfrac12\,\frac{\rho c_E}{T_0}\,\bth^2 .
  \label{eq:psi}
\end{align}

\begin{remark}[Why the shear term carries $\tfrac12 C_{44}\gamma_{xz}^2$ and not $2C_{44}\eps_{xz}^2$]
Because $\gamma_{xz}=2\eps_{xz}$, writing $\tfrac12 C_{44}\gamma_{xz}^2$ is the
same as $2C_{44}\eps_{xz}^2$, which is the correct Voigt form
$\tfrac12\s_{xz}\gamma_{xz}$. Using $\tfrac12 C_{44}\eps_{xz}^2$ (without the
factor) would under-count the shear energy by a factor of two and give
$\s_{xz}=\tfrac12 C_{44}\gamma_{xz}$. This is a common slip; check B1 catches it
(a first version of the check did contain exactly this error and failed).
\end{remark}

\subsection{Constitutive relations}
\label{sec:constitutive}

The constitutive relations are the gradients of $\Psi$. \textsc{Derived}: with
$\s_{ij} = \partial\Psi/\partial\eps_{ij}$,
$\s_{xz} = \partial\Psi/\partial\gamma_{xz}$,
$D_i = -\partial\Psi/\partial E_i$ and $\rho\eta = -\partial\Psi/\partial\bth$,
\begin{align}
  \s_{xx} &= C_{11}\eps_{xx} + C_{12}\eps_{yy} + C_{13}\eps_{zz}
            - e_{31}E_z - \beta_1\bth, \label{eq:sxx}\\
  \s_{yy} &= C_{12}\eps_{xx} + C_{11}\eps_{yy} + C_{13}\eps_{zz}
            - e_{31}E_z - \beta_1\bth, \label{eq:syy}\\
  \s_{zz} &= C_{13}(\eps_{xx}+\eps_{yy}) + C_{33}\eps_{zz}
            - e_{33}E_z - \beta_3\bth, \label{eq:szz}\\
  \s_{xz} &= C_{44}\gamma_{xz} - e_{15}E_x, \label{eq:sxz}\\
  D_x &= e_{15}\gamma_{xz} + \eps_{11}E_x, \label{eq:Dx}\\
  D_z &= e_{31}(\eps_{xx}+\eps_{yy}) + e_{33}\eps_{zz} + \eps_{33}E_z + p\,\bth,
  \label{eq:Dz}\\
  \rho\eta &= \beta_1(\eps_{xx}+\eps_{yy}) + \beta_3\eps_{zz} + p\,E_z
             + \frac{\rho c_E}{T_0}\bth .
  \label{eq:entropy}
\end{align}
Equations~\eqref{eq:sxx}--\eqref{eq:Dz} are exactly Blueprint Eqs.~(2)--(3) with
the sign convention $\s_{ij}=C_{ijkl}\eps_{kl}-e_{kij}E_k-\beta_{ij}\bth$,
$D_i=e_{ikl}\eps_{kl}+\eps_{ik}E_k+p_i\bth$. \textsc{Tested} (check B1).

\begin{remark}[Sign of the pyroelectric term --- a substantive point]
\label{rem:pyro}
In \eqref{eq:Dz} the pyroelectric term is $+p\bth$, while in
\eqref{eq:entropy} it is $+pE_z$ with the \emph{same} sign of $p$, because both
come from the single term $-p\bth E_z$ in \eqref{eq:psi}. Consistency then
requires
\begin{equation}
  \frac{\partial \s_{zz}}{\partial \bth} = -\beta_3,
  \qquad
  \frac{\partial(\rho\eta)}{\partial\eps_{zz}} = \beta_3
  \quad\Longrightarrow\quad
  \frac{\partial^2\Psi}{\partial\bth\,\partial\eps_{zz}} = -\beta_3
  \ \text{both ways,}
  \label{eq:maxwell1}
\end{equation}
and, mixing the electric and thermal variables,
\begin{equation}
  \frac{\partial D_z}{\partial \bth} = p,
  \qquad
  \frac{\partial(\rho\eta)}{\partial E_z} = p
  \quad\Longrightarrow\quad
  \frac{\partial^2\Psi}{\partial\bth\,\partial E_z} = -p \ \text{both ways.}
  \label{eq:maxwell2}
\end{equation}
If the sign of $p$ in \eqref{eq:entropy} were flipped, \eqref{eq:maxwell2} would
fail and no single $\Psi$ would exist. A formulation that inserts a pyroelectric
term into Gauss' law with one sign and into the entropy with the other cannot be
derived from an energy and is thermodynamically inadmissible.
\textsc{Tested} (check B1 evaluates both residuals symbolically; both are zero).
\end{remark}

\subsection{Isotropic limit}
\label{sec:isotropic}

For an elastically isotropic, centrosymmetric medium,
\begin{equation}
  C_{33}=C_{11}=\lambda+2\mu,\quad C_{12}=C_{13}=\lambda,\quad C_{44}=\mu,
  \quad e_{31}=e_{33}=e_{15}=0,\quad p=0,
  \label{eq:isotropic-limits}
\end{equation}
and $C_{13}=C_{11}-2C_{44}$ holds \emph{exactly}. This is used in check I1.
Note that \textbf{PZT-5H does not satisfy this identity} ($C_{13}=84.1$ GPa
against $C_{11}-2C_{44}=79.4$ GPa): it is genuinely transversely isotropic, and an
isotropic closed form must not be applied to it. This distinction is recorded
because it is easy to get wrong, and an earlier version of check I1 did.

\section{Momentum balance and Gauss law}
\label{sec:balance}

With no body force and no free volume charge, and in the quasi-static electric
approximation (the electric field is much slower than $c$; magnetic effects are
neglected),
\begin{equation}
  \s_{ij,j} = \rho\,\ddot u_i, \qquad D_{i,i} = 0 .
  \label{eq:balance}
\end{equation}
These are Blueprint Eq.~(4). \textsc{Assumed} at the level of the quasi-static
approximation; electromagnetically this is standard for the frequency range
considered here and is not re-derived in this note.

In mode form \eqref{eq:balance} reads
\begin{equation}
  i\kappa\,\s_{xx} + i\xi\,\s_{xz} + \rho\omega^2 U = 0, \qquad
  i\kappa\,\s_{xz} + i\xi\,\s_{zz} + \rho\omega^2 W = 0, \qquad
  i\kappa D_x + i\xi D_z = 0 .
  \label{eq:balance-mode}
\end{equation}

\section{Thermal equation: entropy balance and Lord--Shulman}
\label{sec:ls}

This is the step the blueprint flags as requiring independent verification
(Sec.~3 and gate G4), so it is done slowly.

\subsection{Step 1: entropy balance}
The linearised local entropy balance about $T_0$, with no external heat supply, is
\begin{equation}
  T_0\,\frac{\partial(\rho\eta)}{\partial t} = -\nabla\cdot\bq,
  \label{eq:entropy-balance}
\end{equation}
where $\bq$ is the heat flux. (\textsc{Derived} from the Clausius--Duhem
inequality for a linear theory; the factor $T_0$ is what converts entropy
production into power per unit volume.)

\subsection{Step 2: the flux law}
Lord--Shulman replaces Fourier conduction by a single-relaxation (Cattaneo) law,
\begin{equation}
  \bq + \tau_q\,\frac{\partial\bq}{\partial t} = -\kappa\,\nabla\bth,
  \qquad \tau_q \ge 0 .
  \label{eq:cattaneo}
\end{equation}
This is Blueprint Eq.~(5b). \textsc{Assumed} as the thermal model; $\tau_q$ is
\emph{not} an independently measured constant for the materials used (see
\texttt{literature/material\_constants.csv}).

\subsection{Step 3: eliminate the flux}
Apply the operator $(1+\tau_q\partial_t)$ to the divergence of $\bq$. By
\eqref{eq:cattaneo}, $(1+\tau_q\partial_t)(\nabla\cdot\bq) = -\kappa\nabla^2\bth$.
On the other hand, applying the same operator to \eqref{eq:entropy-balance} gives
$(1+\tau_q\partial_t)(\nabla\cdot\bq) = -T_0(1+\tau_q\partial_t)\,\partial_t(\rho\eta)$.
Equating the two,
\begin{equation}
  \kappa_{ij}\,\bth_{,ij}
  = (1+\tau_q\partial_t)\Big[\rho c_E\,\dot\bth
    + T_0\big(\beta_{ij}\dot\eps_{ij} + p_i\dot E_i\big)\Big],
  \label{eq:ls}
\end{equation}
using \eqref{eq:entropy} for $\rho\eta$. \textsc{Derived}, and \textsc{Tested}:
check B2 performs the elimination symbolically and finds a residual of exactly
zero.

\begin{remark}[The operator multiplies the whole bracket]
\label{rem:bracket}
The $(1+\tau_q\partial_t)$ factor cannot be attached to the mechanical term alone.
It arises because the flux law \eqref{eq:cattaneo} is a statement about the
\emph{total} heat flux, and that flux also carries the thermoelastic and
pyroelectric sources through \eqref{eq:entropy-balance}. Putting the relaxation
factor only on $\rho c_E\dot\bth$ (a form that does appear in some papers) is a
different theory, not a rearrangement. This is the ``sign and entropy
verification'' that Blueprint V5.0 asks for, and the answer is that the blueprint's
Eq.~(5) is correct.
\end{remark}

\subsection{Step 4: time-harmonic form}
With $e^{-i\omega t}$, $\partial_t\to -i\omega$, so
\begin{equation}
  (1+\tau_q\partial_t) \longrightarrow (1 - i\omega\tau_q) \equiv \Lambda .
  \label{eq:Lambda}
\end{equation}
The mode form of \eqref{eq:ls} is
\begin{equation}
  -\big(\kappa_{11}\kappa^2 + \kappa_{33}\xi^2\big)\Theta
  + i\omega\Lambda\Big[\rho c_E\,\Theta
    + T_0\big(\beta_1\eps_{xx} + \beta_3\eps_{zz} + p\,E_z\big)\Big] = 0 .
  \label{eq:ls-mode}
\end{equation}
Note the sign of $\Lambda$: with $e^{-i\omega t}$ the relaxation factor is
$1-i\omega\tau_q$, whose imaginary part is negative. Using $1+i\omega\tau_q$
would correspond to $e^{+i\omega t}$ and, if mixed with the other convention,
produces a growing thermal mode. This is a loss-sign trap and is recorded in
\texttt{reports/UNRESOLVED\_ISSUES.md} (N-1).

\section{Stability and admissibility conditions}
\label{sec:stability}

The quadratic form \eqref{eq:psi} must be positive definite for the reference
state to be stable. In the decoupled limits this requires
\begin{equation}
  C^{E} \succ 0,\qquad \eps^{S} \succ 0,\qquad
  \rho>0,\quad c_E>0,\quad T_0>0,\quad \kappa_{ij}\succ 0,\quad \tau_q\ge 0 .
  \label{eq:stability}
\end{equation}
\textsc{Tested} (check A1 for $C$ and $\eps$, check H1 for the thermal set,
for all four materials in \texttt{literature/material\_constants.csv}).

\paragraph{Causality floor on $\tau_q$.} The Lord--Shulman theory admits a
thermal wave of speed $v_T = \sqrt{\kappa/(\rho c_E\tau_q)}$. Requiring it not to
exceed the elastic longitudinal speed $c_L$ gives
\begin{equation}
  \boxed{\ \tau_q \ \ge\ \tau_{q,\min} \;=\; \frac{\kappa}{\rho\,c_E\,c_L^{2}}\ }
  \label{eq:tau-floor}
\end{equation}
\textsc{Derived}. Values: PZT-5H $2.83\times10^{-14}$ s, BaTiO$_3$
$1.45\times10^{-14}$ s, Al6061 $1.73\times10^{-12}$ s, CdSe
$4.67\times10^{-13}$ s. Every working value in the material table is chosen
\emph{above} the floor and the ratio $\tau_q/\tau_{q,\min}$ is reported by check
H1. This is a constraint on admissible parameter ranges, not a fit; no value was
adjusted to produce a physical-looking result.

\section{Dimensional audit}
\label{sec:dims}

Written in the exponent order $(\mathrm{kg},\mathrm{m},\mathrm{s},\mathrm{A},\mathrm{K})$:
\begin{center}
\begin{tabular}{lll}
\toprule
Quantity & Dimension & SI \\
\midrule
$C_{ijkl}$ & $(1,-1,-2,0,0)$ & Pa \\
$e_{kij}$ & $(0,-2,1,1,0)$ & C\,m$^{-2}$ \\
$\eps_{ij}$ & $(-1,-3,4,2,0)$ & F\,m$^{-1}$ \\
$\beta_{ij}$ & $(1,-1,-2,0,-1)$ & Pa\,K$^{-1}$ \\
$p_i$ & $(0,-2,1,1,-1)$ & C\,m$^{-2}$K$^{-1}$ \\
$\rho$ & $(1,-3,0,0,0)$ & kg\,m$^{-3}$ \\
$c_E$ & $(0,2,-2,0,-1)$ & J\,kg$^{-1}$K$^{-1}$ \\
$\kappa_{ij}$ & $(1,1,-3,0,-1)$ & W\,m$^{-1}$K$^{-1}$ \\
$u_i$ & $(0,1,0,0,0)$ & m \\
$\varphi$ & $(1,2,-3,-1,0)$ & V \\
$\bth$ & $(0,0,0,0,1)$ & K \\
\bottomrule
\end{tabular}
\end{center}

The four rows of the characteristic matrix then carry the dimensions
$\mathrm{N\,m^{-3}}$, $\mathrm{N\,m^{-3}}$, $\mathrm{C\,m^{-3}}$ and
$\mathrm{W\,m^{-3}}$ respectively. \textsc{Tested}: check A2 verifies all 27
independent entries of $M$ against these targets and finds zero mismatches.

\begin{remark}[What was \emph{not} a valid test]
An earlier version of check A2 tested homogeneity of $M$ under
$(\kappa,\xi)\mapsto(\kappa/L,\xi/L)$. That test is invalid: $\omega$ is an
inverse time and does not rescale with length, so the terms $\omega^2\rho$ and
$c_E\omega\rho$ necessarily break such a scaling. Dimensional analysis is the
correct test. The failed version is recorded rather than deleted because the
distinction matters for anyone extending the code.
\end{remark}

\begin{remark}[What the thermal coupling does \emph{not} do]
\label{rem:sh}
For $6mm$ symmetry with poling along $z$, $\beta_{ij}=\mathrm{diag}(\beta_1,\beta_1,\beta_3)$
and $p_i=(0,0,p)$. Anti-plane motion $u_y(x,z)$ with $\eps_{xx}=\eps_{zz}=\eps_{yy}=0$
and $E_y=E_y(x,z)$ gives $\beta_{ij}\eps_{ij}=0$ and $p_iE_i = 0$ identically.
\textbf{The thermal field therefore decouples from SH motion exactly.} This is a
structural fact about $6mm$ symmetry, not a numerical accident and not a
modelling choice, and it means SH configurations cannot be used to demonstrate any
thermal-coupling effect. \textsc{Derived}.
\end{remark}

\section{Summary of what is established, and what is not}
\label{sec:summary}

\paragraph{Established (\textsc{Derived} + \textsc{Tested}).}
Equations~\eqref{eq:strain-mode}--\eqref{eq:balance-mode} and \eqref{eq:ls-mode}
constitute a closed, thermodynamically admissible bulk system, generated by the
single potential \eqref{eq:psi}. All constitutive relations, the entropy relation
and the Lord--Shulman equation reproduce the blueprint's Eqs.~(2)--(5) exactly,
including the placement of the relaxation operator.

\paragraph{Not established.}
\begin{itemize}
  \item The pyroelectric coefficient $p$ and the thermoelastic couplings
        $\beta_1,\beta_3$ are \textsc{Assumed} order-of-magnitude values for
        PZT-5H and BaTiO$_3$; only the CdSe set has tabulated $\beta$ and $p$.
        No quantitative result may rest on their precise values.
  \item $\tau_q$ is a stated choice above a derived causal floor, not a
        measurement.
  \item The quasi-static electric approximation is assumed, not re-derived.
  \item No wave-propagation, reflection or transmission result appears in this
        note. Those are in \texttt{characteristic\_matrix.tex} and
        \texttt{energy\_flux\_derivation.tex}.
\end{itemize}

\end{document}
